% ============================================================================
\section{Machine-checked formalization}\label{sec:lean}
% ============================================================================
The obstruction is short but easy to misstate: data entry, terminal attitude, state constraints and the sign conventions of
the dive-then-brake face all matter. We therefore formalized its core in Lean~4 with Mathlib \cite{demoura2021lean4,mathlib2020}
(toolchain \texttt{v4.36.0-rc2}; project \texttt{lean/LcvxObstruction}, namespace \texttt{LcvxObstruction}). The project contains
no \texttt{sorry} and no \texttt{axiom}. For every headline theorem, \texttt{\#print axioms} reports exactly
\[
\texttt{[propext, Classical.choice, Quot.sound]},
\]
the standard axioms of classical Lean mathematics. This was checked after a full \texttt{lake build} at commit \texttt{eb2f850}.

\subsection{What is formalized}
\Cref{tab:lean} lists the formal theorems that correspond to statements in this paper. A row-by-row faithfulness map from each
Lean definition to its mathematical meaning is maintained in \texttt{lean/MODEL.md}. Its main conventions are listed next.
\begin{itemize}[leftmargin=*,itemsep=2pt]
\item $\R^3$ is \lean{EuclideanSpace \ensuremath{\mathbb{R}} (Fin 3)}. Angles are Mathlib's \lean{InnerProductGeometry.angle}. Integrals are Bochner
integrals over $[0,t_f]$.
\item \lean{UniformConvexRelaxation \ensuremath{\Theta} Y} is a convex set $C\subseteq\Theta\times Y$ for arbitrary real vector spaces $\Theta,Y$.
\lean{IsLossless F} is the feasibility form of \cref{def:lossless}(a), and \lean{IsLosslessFor J F} is the optimization (LCvx) form
(b) for an arbitrary cost $J$, including attainment.
\item \lean{AttitudeAdmissible tf \ensuremath{\omega} \ensuremath{\delta} d b} encodes M1 with variant (F): $b(0)=\ez$, unit vectors,
$\angle(b(s),b(s'))\le\omega|s-s'|$, measurable unit $d$, gimbal $\angle(d,b)\le\delta$. Pointwise constraints are imposed for every $s$,
which changes no integral. The admissible thrust accelerations form an abstract class \lean{Acc}. Theorem B assumes only
$a\ge a_1$ for a continuous $a_1>0$ (so it covers mass depletion through $a_1(s)=\rho_1/(m_0-\alpha\rho_1s)$). Theorem B$^*$ assumes a
box $a_{\min}\le a\le a_{\max}$ with constant bounds, which is exact for constant mass.
\item Positions and velocities are given in closed (Cauchy) form, $r(t)=r_0+tv_0+\tfrac{t^2}2g+\int_0^t(t-s)\,a\,d$.
\item \lean{StateLinearRelaxation} is \cref{def:statelinear} with an affine readout and a convex penalty that is nonnegative on $C$.
\end{itemize}

\begin{table}[t]\centering\footnotesize
\caption{Lean theorems corresponding to results in this paper. All compile without \texttt{sorry}; axioms
\texttt{propext}, \texttt{Classical.choice}, \texttt{Quot.sound} only.}\label{tab:lean}
\begin{tabular}{@{}>{\raggedright\arraybackslash}p{0.22\linewidth}>{\raggedright\arraybackslash}p{0.44\linewidth}>{\raggedright\arraybackslash}p{0.28\linewidth}@{}}\toprule
Paper & Lean name (file) & Formal scope \\\midrule
\Cref{thm:A}, \cref{prop:opt-feas} & \lean{theoremA}, \lean{not\_isLossless\_of\_not\_convex}, \lean{not\_isLosslessFor\_of\_not\_convex}, \lean{IsLosslessFor.isLossless} (\texttt{Abstract}) & any $\Theta,Y$, any cost \\
Spherical triangle inequality and rigidity & \lean{inner\_triangle}, \lean{inner\_triangle\_eq\_lateral}, \lean{lateral\_eq\_of\_inner\_ge} (\texttt{Sphere}) & inner-product form with equality case \\
\Cref{thm:B} (variant F) & \lean{not\_convex\_reachSet}, \lean{not\_convex\_landingSet}, \lean{no\_lossless\_relaxation\_5dof}, \lean{no\_losslessFor\_relaxation\_5dof}, \lean{landingSet\_nonempty} (\texttt{Nonconvex}, \texttt{Main}) & $\delta+\omega t_f\le\pi/2$, any continuous $a_1>0$, any $g$, no state constraints, $b(t_f)$ free \\
\Cref{thm:Bstar} (qualitative) & \lean{VStar\_not\_mem\_WSet}, \lean{witness\_mem\_glideSlope}, \lean{not\_convex\_constrainedLandingSet}, \lean{exists\_switchTime}, \lean{no\_lossless\_relaxation\_constrained\_5dof}, \lean{no\_losslessFor\_relaxation\_constrained\_5dof}, \lean{constrainedLandingSet\_nonempty} (\texttt{DiveBrake}, \texttt{Main}) & constant acceleration bounds ($\alpha=0$), $g_0<a_{\max}$, every $t_f>0$, every $\X\supseteq$ glide-slope set, $b(t_f)$ free \\
\Cref{thm:D} (abstract) & \lean{StateLinearRelaxation.convexHull\_subset}, \lean{minError\_lossy}, \lean{convex\_superset\_gap} (\texttt{TheoremD}) & abstract; concrete instance not formalized \\
\Cref{lem:LS} (abstract) & \lean{not\_lineSearch\_lossless}, \lean{symmetric\_supply}, \lean{not\_lineSearch\_lossless\_of\_symmetry} (\texttt{FreeTime}) & abstract (i)--(ii); sorry-free (commit \texttt{517f1bb}); axiom check after next full build \\\bottomrule
\end{tabular}
\end{table}

\subsection{The PMP-as-hypothesis convention}
Two external theorems appear on the positive side, and each enters only as an explicit hypothesis of the formal statement,
never as an axiom. The Pontryagin maximum principle is a structure \lean{PMPConclusion} whose fields are the adjoint equations
and the a.e.\ minimum condition of (H-PMP). Aumann's theorem enters as \lean{AumannInstance}, the single instance used in
\cref{thm:CW}. A reader can therefore check exactly which conclusions depend on which external results. The positive-side files
(\texttt{Positive/}: pointwise cap lemmas, primer geometry, \cref{thm:CS} under \lean{PMPConclusion}, discrete KKT and purification,
\cref{thm:C6} pointwise and trajectory-level convexity, \cref{thm:CW} feasibility under \lean{AumannInstance}) are summarized in
\texttt{lean/POSITIVE.md}. There the Lean form is weaker than the paper in several places, which that file lists explicitly: the
adjoint ODEs are assumed globally differentiable; C-S(2) is constant-mass only; ``$\mathrm P_0$-optimal'' is reduced to pointwise
feasibility; C-E is discrete-time; and C-W is constant-mass and unconstrained.

\begin{todobox}[{(lean, lean-s, lean-positive): formalizations still pending}]
\textbf{Not yet committed or not yet axiom-audited at the time of writing:}
\begin{itemize}[leftmargin=*,itemsep=1pt]
\item \texttt{SixDOF.lean} (\cref{thm:6kin}): \lean{not\_convex\_landingSet6}, \lean{not\_convex\_constrainedLandingSet6}. Status per [lean]: PENDING.
\item \texttt{Slice/} (lean-s): \lean{periodic\_wirtinger}, \lean{wirtinger\_scalar}, \lean{wirtinger\_lipschitz} (\cref{lem:W}); \lean{master\_inequality},
\lean{dent} (master inequality of \cref{thm:S}, constant mass); \lean{mixing}, \lean{aOpt}, \lean{no\_valueLossless\_of\_strict} (\cref{thm:Aopt}).
Lemma S itself (the witness and $0\notin\cl L$) and Theorem B$^*$ beyond the qualitative constant-mass form are not formalized.
\item \texttt{Positive/} (lean-positive): built locally, but the axiom audit output in \texttt{lean/POSITIVE.md} is a placeholder, and the files are not
yet in the root import list.
\end{itemize}
\textbf{To fill in:} move each item into \cref{tab:lean} once it is committed with a clean \texttt{\#print axioms}.
\end{todobox}

\subsection{What is not formalized}
Mass depletion in Theorem B$^*$; the terminal-attitude variants (T) and (T$'$); free final time for B$^*$ (the abstract line-search
lemma is formalized, the supply of its hypotheses from B$^*$ is not); all quantitative gaps (\cref{thm:Q}, the B$^*$ bound,
$\varepsilon$-losslessness); 6-DOF and torque models; the centre-of-percussion identity; $\rho_1=0$; the SDP results; the equivalence of
the closed-form trajectory with $r_0+\int v$ (standard Fubini). The physical scope caveat of \cref{sec:Bstar} also applies to the
formal Theorem B: its witnesses start below ground for a vehicle that cannot hover at minimum throttle. The formal Theorem B$^*$ is
the statement that covers the altitude-constrained problem.
